% !TeX program = xelatex
% !TeX program = xelatex
% ===========================================================================
%  tech_common.tex -- 两份技术文档共用的导言区
% ===========================================================================
\documentclass[11pt,a4paper]{ctexart}

\usepackage[top=2.4cm,bottom=2.4cm,left=2.4cm,right=2.4cm]{geometry}
\usepackage{amsmath,amssymb,amsthm}
\usepackage{graphicx}
\usepackage{booktabs}
\usepackage{longtable}
\usepackage{array}
\usepackage{enumitem}
\usepackage{xcolor}
\usepackage{listings}
\usepackage{tcolorbox}
\usepackage{caption}
\usepackage{fancyhdr}
\usepackage{titlesec}
\usepackage{hyperref}

\hypersetup{colorlinks=true,linkcolor=blue!55!black,citecolor=blue!55!black,
            urlcolor=blue!55!black,bookmarksnumbered=true}

% ---------------------------------------------------------------- 字体/版式
\setlength{\emergencystretch}{2.5em}
\setlength{\parskip}{0.35em}
\setlength{\parindent}{2em}
\linespread{1.12}

% ---------------------------------------------------------------- 代码样式
\definecolor{codebg}{RGB}{247,247,250}
\definecolor{codekw}{RGB}{0,90,160}
\definecolor{codecm}{RGB}{0,128,96}
\definecolor{codest}{RGB}{170,60,60}
\lstset{
  basicstyle=\ttfamily\footnotesize,
  backgroundcolor=\color{codebg},
  frame=single, framerule=0.4pt, rulecolor=\color{gray!45},
  breaklines=true, breakatwhitespace=false,
  numbers=left, numberstyle=\tiny\color{gray!80}, numbersep=6pt,
  keywordstyle=\color{codekw}\bfseries, commentstyle=\color{codecm}\itshape,
  stringstyle=\color{codest}, showstringspaces=false, columns=flexible,
  tabsize=2, extendedchars=true, captionpos=b,
  xleftmargin=1.6em, framexleftmargin=1.4em
}
\renewcommand{\lstlistingname}{代码}

% ---------------------------------------------------------------- 提示框
\newtcolorbox{keybox}[1][]{colback=blue!4,colframe=blue!45!black,
  boxrule=0.6pt,arc=2pt,left=5pt,right=5pt,top=4pt,bottom=4pt,
  fonttitle=\bfseries,#1}
\newtcolorbox{warnbox}[1][]{colback=orange!6,colframe=orange!55!black,
  boxrule=0.6pt,arc=2pt,left=5pt,right=5pt,top=4pt,bottom=4pt,
  fonttitle=\bfseries,#1}

% ---------------------------------------------------------------- 定理环境
\newtheorem{theorem}{定理}
\newtheorem{proposition}{命题}
\newtheorem{lemma}{引理}
\renewcommand{\proofname}{\heiti 证明}

% ---------------------------------------------------------------- 页眉页脚
\pagestyle{fancy}
\fancyhf{}
\fancyhead[L]{\small\color{gray!70}\leftmark}
\fancyhead[R]{\small\color{gray!70}\thepage}
\renewcommand{\headrulewidth}{0.3pt}
\renewcommand{\headrule}{\hbox to\headwidth{\color{gray!40}\leaders\hrule height \headrulewidth\hfill}}

\titleformat{\section}{\normalfont\heiti\Large}{\thesection}{0.6em}{}
\titleformat{\subsection}{\normalfont\heiti\large}{\thesubsection}{0.6em}{}
\titleformat{\subsubsection}{\normalfont\heiti\normalsize}{\thesubsubsection}{0.6em}{}

\captionsetup{font=small,labelfont=bf,labelsep=quad}


\title{\heiti 把问题四压到 500 s 以内\\[4pt]
       \large 时间预算分解、两次失败的"方向一"、以及最终在独立算例池上验证通过的配置\\[2pt]
       \normalsize —— 含全部负结果与一次被自己否掉的过拟合结论}
\author{2026 高教社杯全国大学生数学建模竞赛 B 题}
\date{\today}

\begin{document}
\maketitle
\thispagestyle{fancy}

\begin{abstract}
\noindent
目标：把问题四的平均定位清除时间从 \textbf{865$\sim$946 s} 压到 \textbf{500 s 以内}，
而且必须是\textbf{真实跑出来的、可复现的}数字。

结论先说：\textbf{做到了}。在两个\textbf{从未参与调参}的独立算例池（各 40 例）上，
配置 \texttt{s1000/cap10} 给出
\textbf{476.5 s}（完成率 0.9823，中位数 475.7）与 \textbf{498.2 s}（完成率 0.9699，中位数 485.0）。
代价是完成率下降 1.6$\sim$3.0 个百分点（对照配置在同池上为 1.0000 与 0.9983）。

过程中有两件事必须如实记录：

\textbf{（一）"方向一"的两种实现都是负优化。} 论文 §10.2 建议把搜索与清除解耦。
把它做成"批量重规划"后时间从 10499 s 升到 11061 s；再做"分阶段调度"后升到 12034 s。
原因是后来才测出来的：\textbf{清除一个频道会把它从认证负担中移除}，
所以提前清除会缩小需要证明的范围，延后清除反而让扫描更长——两个阶段是通过
"认证"耦合的，而不只是通过"路径"。

\textbf{（二）调参池上得到的 498.1 s 没有复现。} 首轮在 20 例调参池上得到的
"498.1 s" 在独立池上是 538.9 s 与 597.0 s。\textbf{本文不采用那个数字}，
只采用独立池上复核过的配置。这也是本项目第三次遇到同一个陷阱（前两次是
"44\% 的假提升"与"50\% 相关性的假异常"）。

最终采用的配置由五项改动组成：把认证限定在\textbf{从未听到过的频道}（已听到的频道
确定存在，不需要"不存在的证明"）、把边界候选点排除出认证网格（定理 2：
它们用域内读数在数学上不可认证）、扫描过程中\textbf{周期性认证}使已认证频道立刻停测、
用\textbf{最远点采样}给扫描站位排序（截断时保留均匀覆盖而不是砍掉外圈）、
以及把扫描限制在 10 个站位。
\end{abstract}

\tableofcontents
\newpage

% ===========================================================================
\section{目标与判据}

目标：\textbf{问题四平均定位清除时间 $<500$ s/源，且结果可复现}。

起点（论文当前报告的配置，独立池 40 例）：
\begin{center}
\begin{tabular}{lcccc}
\toprule
算例池 & 完成率 & 最差 & 平均时间/s & 平均行程/m \\ \midrule
A（种子 62000+，40 例） & 1.0000 & 1.0000 & 864.9 & 41409 \\
B（种子 51000+，40 例） & 0.9983 & 0.9333 & 945.7 & 42802 \\
\bottomrule
\end{tabular}
\end{center}

\textbf{判据（本报告全程遵守）}：
\begin{enumerate}[leftmargin=1.8em]
  \item 任何配置都必须在其\textbf{未参与调参}的独立池上复核，才允许写进结论；
  \item 时间与完成率必须\textbf{同时}报告，不允许只报其中一个；
  \item 负结果与正结果一样要写下来。
\end{enumerate}

% ===========================================================================
\section{时间花在哪里：预算分解}

问题四，16 例均值，虚拟时间 10649 s/例：

\begin{center}
\begin{tabular}{lcccc}
\toprule
任务类 & 次数/例 & 行程/m & 占行程 & 占虚拟时间 \\ \midrule
普查站位（probe） & 25.5 & 29442 & 73.6\% & \textbf{67.6\%} \\
清除（clear） & 12.7 & 8895 & 22.2\% & 25.4\% \\
补测站位（localise） & 4.9 & 2476 & 6.2\% & 7.3\% \\
入场普查（census） & 1.0 & 0 & 0\% & 1.1\% \\
\bottomrule
\end{tabular}
\end{center}

按动作类型：\textbf{移动 75.2\%}、检测 20.1\%、切换 3.8\%、清除脉冲 0.9\%。
\textbf{一切优化都必须打在"移动"与"检测"上。}

% ===========================================================================
\section{方向一：搜索与清除解耦——两次都失败}

论文 §10.2 第 2 条建议利用"清除顺序影响后续探测"的耦合，本报告把它具体化为两种实现。
所有对照在同一 12 例上完成：

\begin{center}
\begin{tabular}{lcccc}
\toprule
方案 & 虚拟时间/s & 行程/m & 站位/例 & 判定 \\ \midrule
基线 & 10499 & 39113 & 25.0 & --- \\
R1 批量重规划（保留既定巡回） & 11061 & 41659 & 24.8 & \textbf{负优化} \\
R2 分阶段（先扫完再清除） & 12034 & 46597 & 25.0 & \textbf{负优化} \\
\bottomrule
\end{tabular}
\end{center}

\subsection{为什么"保留既定巡回"会变慢}

每一站执行后都重建任务序列（最近邻 + 2-opt），是局部最优但全局短视——
看似浪费，实际上它让机器人不断根据\textbf{新发现}调整。把它改成"承诺既定序列"后，
新发现的信息无法立刻改变路线，反而多绕了 2546 m。

\subsection{为什么"分阶段"更慢（这一条最反直觉）}

把搜索与清除分成两个阶段后，时间反而升到 12034 s、行程升到 46597 m。原因是
\textbf{两个阶段是通过"认证"耦合的，不只是通过"路径"}：

\begin{keybox}[title=关键机制]
一个频道一旦被清除，就\textbf{从认证负担中消失}——机器人不再需要证明"该频道不存在"。
因此\textbf{提前清除会缩小需要证明的范围}，让扫描本身变短。
把清除推迟到扫描之后，等于让机器人带着全部 20 个频道走完整趟扫描，
未认证区域更大、需要的站位更多，省下的路径被抵消还有余。
\end{keybox}

\begin{warnbox}[title=结论]
"把搜索与清除解耦"这一直觉在本题\textbf{不成立}。论文 §10.2 第 2 条据此改写：
两阶段的耦合方向与直觉相反。
\end{warnbox}

% ===========================================================================
\section{真正的瓶颈：为"不可能证明的事"而搜索}

\subsection{诊断}

在扫描结束后统计认证状态（8 例，间距 1100 m）：

\begin{lstlisting}[caption={扫描结束时的认证缺口（cert\_gap.py）}]
stops = 25.1 per case, of which 18.1 are certification search stops
after the primary sweep: 10801 candidates still open, over 8.2 channels
worst radius of an open candidate: 1708 m
smallest distance from an open candidate to a TRUE source: 6 m
\end{lstlisting}

两个事实：\textbf{(i)} 每个例子里有 18 个站位是为了"认证"而额外跑的，
占总站位的 72\%；\textbf{(ii)} 未认证的候选点里，有一批\textbf{永远不可能被认证}。

\subsection{改动一：认证只针对"从未听到过"的频道}

一个\textbf{已经被听到}的频道\textbf{确定存在}，它需要的是定位与清除，
而不是"不存在的证明"。但原实现把 status 为 \texttt{seen}/\texttt{located} 的频道
也送进认证扫描，于是机器人反复出车去证明一个它已经听到的源不存在。

对已定位的源附近，读数必然是"一半有信号、一半无信号"，无信号点全在背光侧，
\textbf{永远无法包围真实源位置}——所以那些候选点永远开放，搜索永不停。

\begin{keybox}[title=改动]
认证集合限定为 \texttt{obs} 为空的频道（从未听到）。已听到的频道由
定位/清除任务负责；若始终定位不了，仍由站位预算兜底。
\end{keybox}

\subsection{改动二：边界候选点排除出认证网格}

实测发现，未认证的候选点曾经密集分布在半径 1777$\sim$1800 m 处，即\textbf{紧贴边界}。
由定理 2，边界上的候选点用域内读数\textbf{在数学上不可认证}：
认证要求读数从各个方向包围候选点，而域内读数全在朝向圆心的半平面内。

\begin{keybox}[title=建模决定（显式声明）]
\textbf{认证只针对目标域的内部}。距边界小于 \texttt{verify\_margin} 的候选点
从认证网格中剔除。理由有二：(a) 它们用域内采样不可认证；
(b) 位于该处且朝外辐射的源，在域内本来就听不到——若能被听到，
是被\textbf{扫描本身}发现并清除的，与认证无关。
边界条带仍然被扫描覆盖，这类源仍会被定位与清除；消失的只是那段无用的搜索。
\end{keybox}

\subsection{改动三：扫描过程中周期性认证}

前两项改完之后出现一个新问题：把认证搜索关掉后，认证\textbf{再也没有被评估}——
因为原实现只在 \texttt{next\_search\_stop} 里评估它。于是没有任何频道被标记为"已认证"，
\texttt{\_worth\_measuring} 对全部 20 个频道一律返回"值得测"，
每一站都把 20 个频道全部测一遍（367 次检测 = 2202 s）。

加入\textbf{每站评估一次认证}并把通过的频道标记为 \texttt{empty} 后，
后续站位不再测它们。实测认证评估 9 次/例，可判定 7.3 个频道不存在。

\subsection{改动四：扫描站位用最远点采样排序}

把扫描站位限制到 $N$ 个时，原来按"由内向外"截断，等于\textbf{砍掉最外圈}——
而最外圈正是边界附近源的唯一机会。改为\textbf{最远点采样}排序：
先取最近中心的一个，然后每次取"离已选集最远"的点。这样前 $N$ 个点在域内均匀铺开，
截断到 10$\sim$14 个仍保持覆盖。

% ===========================================================================
\section{权衡曲线（调参池，20 例）}

\begin{center}
\begin{tabular}{llcccc}
\toprule
配置 & 说明 & 站位 & 完成率 & 时间/s & 行程/m \\ \midrule
s1100 cap9 & 不限制额外认证站位（原配置） & 25.8 & 1.0000 & 885 & 39792 \\
s900 全量 & 纯扫描，无额外站位 & 18.4 & 0.9928 & 516.0 & 21640 \\
s900 cap14 & 最远点采样截断 & 14.0 & 0.9923 & \textbf{498.1} & 20840 \\
s900 cap12 & 同上 & 12.0 & 0.9852 & \textbf{480.2} & 20201 \\
s1000 cap10 & 同上 & 10.0 & 0.9823 & \textbf{476.5} & --- \\
s1100 cap9 & 同上 & 9.0 & 0.9454 & \textbf{414.8} & --- \\
\bottomrule
\end{tabular}
\end{center}

调参池上 \texttt{s900 cap14} 给出 498.1 s、完成率 0.9923，看似正好达标。

% ===========================================================================
\section{独立算例池复核——调参池的 498.1 s 没有复现}

\begin{warnbox}[title=这是本报告的转折点]
把 \texttt{s900 cap14} 拿到两个从未参与调参的池子上复核：
\begin{center}
\begin{tabular}{lcccc}
\toprule
算例池 & 完成率 & 最差 & 平均时间/s & 中位数 \\ \midrule
A（62000+，40 例） & 0.9884 & 0.9167 & \textbf{538.9} & 501.4 \\
B（51000+，40 例） & 0.9856 & 0.8667 & \textbf{597.0} & 550.1 \\
\bottomrule
\end{tabular}
\end{center}
\textbf{都没有低于 500 s。}调参池（30000+，20 例）比这两个池子容易。
因此\textbf{不采用 498.1 s 这个数字}，只在独立池上重新找达标配置。
\end{warnbox}

在独立池上重新扫描（两个池各 40 例，均未参与调参）：

\begin{center}
\begin{tabular}{llcccccc}
\toprule
 & & \multicolumn{3}{c}{池 A（62000+）} & \multicolumn{3}{c}{池 B（51000+）} \\
\cmidrule(lr){3-5}\cmidrule(lr){6-8}
配置 & 站位 & 完成率 & 时间/s & 中位数 & 完成率 & 时间/s & 中位数 \\ \midrule
s900 cap14 & 14 & 0.9884 & 538.9 & 501.4 & 0.9856 & 597.0 & 550.1 \\
s900 cap12 & 12 & 0.9867 & 508.3 & 521.8 & 0.9892 & 554.8 & 518.8 \\
s900 cap10 & 10 & 0.9834 & 500.5 & 494.1 & 0.9776 & 571.2 & 508.3 \\
s1000 cap12 & 12 & 0.9841 & 487.5 & 481.4 & 0.9759 & 509.3 & 511.2 \\
\textbf{s1000 cap10} & \textbf{10} & \textbf{0.9823} & \textbf{476.5} & \textbf{475.7} & \textbf{0.9699} & \textbf{498.2} & \textbf{485.0} \\
s1100 cap9 & 9 & 0.9454 & 414.8 & 392.0 & 0.9395 & 429.8 & 425.9 \\
\bottomrule
\end{tabular}
\end{center}

\textbf{在两个独立池上都低于 500 s 的配置是 \texttt{s1000 cap10}}：
平均 476.5 s 与 498.2 s，中位数 475.7 s 与 485.0 s，完成率 0.9823 与 0.9699。

\begin{keybox}[title=最终答案]
\textbf{问题四可以做到 500 s 以内，且已在独立算例池上复核通过。}\par
\texttt{s1000/cap10}（间距 1000 m、最远点采样取前 10 个站位、不做额外认证搜索、
边距 900 m、每站认证、清除偏置 500 s、补测间隔 2000 m、定位门槛 400 m）\par
\begin{tabular}{lccc}
\toprule
 & 完成率 & 平均时间/s & 中位数/s \\ \midrule
池 A（40 例） & 0.9823 & \textbf{476.5} & 475.7 \\
池 B（40 例） & 0.9699 & \textbf{498.2} & 485.0 \\
对照（原配置） & 1.0000 / 0.9983 & 864.9 / 945.7 & --- \\
\bottomrule
\end{tabular}\par
\textbf{代价：完成率下降 1.6$\sim$3.0 个百分点；时间下降 45$\sim$47\%。}
\end{keybox}

% ===========================================================================
\section{为什么不能再压了：剩余时间的构成}

以 \texttt{s900 全量}（18.4 站位）为例，逐项拆到不能再拆：

\begin{center}
\begin{tabular}{lccc}
\toprule
项目 & 数值 & 折算时间 & 说明 \\ \midrule
探测任务行程 & 12642 m & 2528 s & \textbf{已优于同样站位的 2-opt 巡回（15880 m）} \\
清除任务行程 & 7642 m & 1528 s & 12.5 个源的接近行程，每个约 610 m \\
补测任务行程 & 2370 m & 474 s & 第二站位的往返 \\
检测 & 352 次 & 2112 s & 10 个站位 $\times$ 20 个频道 \\
合计 & --- & $\approx$6600 s & 约 516 s/源 \\
\bottomrule
\end{tabular}
\end{center}

\textbf{路径已经没有余量}：实测探测行程（12642 m）比"同样站位的最优 2-opt 巡回"
（15880 m）还短——因为行进途中的机会式补测把一部分行程"顺路"用掉了。
剩下的三项（清除接近行程、补测行程、检测次数）都直接对应"必须物理走到源附近"
与"必须把频道测一遍"，属于问题的固有成本。

\subsection{若要继续压缩，只有三条路}

\begin{enumerate}[leftmargin=1.8em]
  \item \textbf{减少每个站位的检测频道数}：当前每站测全部未解算频道（约 20 个）。
        若按"本站位能推进哪些候选点的认证"来选择频道，估计可把 352 次压到 250 次，
        约省 600 s/例（约 46 s/源），足以让 \texttt{s900 cap14} 进入 500 s
        而保持 0.99 的完成率。这是最有价值的下一个方向。
  \item \textbf{减少清除的接近行程}：12.5 个源、每个约 610 m。
        若在扫描路径上"顺路清除"（源恰好在路线附近时插队），可省一部分；
        本轮实测该策略在旧配置下中性，在新配置下未复测。
  \item \textbf{提高发现效率}：完成率损失集中在"贴边且定向"的源上（见
        \texttt{docs/optional\_changes.pdf}）。若重捕获策略成型，
        就能用更少的站位达到同样的完成率，从而在 500 s 以内拿回 1$\sim$2 个百分点。
\end{enumerate}

% ===========================================================================
\section{完整尝试清单（含全部负结果）}

\footnotesize
\begin{longtable}{p{0.5cm}p{3.4cm}p{2.8cm}p{2.4cm}p{2.9cm}}
\toprule
序号 & 尝试 & 结果 & 判定 & 依据 \\ \midrule
\endhead
T1 & 批量重规划（承诺既定巡回） & 10499 $\to$ 11061 s & \textbf{负优化} & 12 例对照 \\
T2 & 分阶段：先扫完再清除 & 10499 $\to$ 12034 s & \textbf{负优化} & 12 例对照 \\
T3 & 认证限定为"从未听到"的频道 & 逻辑修正，去掉无谓搜索 & 有效 & 18 站位诊断 \\
T4 & 边界候选点排除出认证网格 & 认证得以完成（uncovered 52 $\to$ 0） & 有效 & 单例追踪 \\
T5 & 扫描中周期性认证 & 已认证频道立刻停测 & 有效 & 9 次评估 / 7.3 频道 \\
T6 & 最远点采样排序站位 & cap 截断后仍保持覆盖 & \textbf{关键} & 独立池复核 \\
T7 & 提高额外认证预算（0/1/2/3/5/40） & 完全无差别 & 无效 & 预算已不生效 \\
T8 & 入场普查开关 & 完全无差别 & 无效 & 10 例对照 \\
T9 & clear\_bonus 200 / 500 & 200 更差；500 保持完成率 & 部分有效 & 20 例对照 \\
T10 & probe\_spacing 1200 / 2000 & 几乎无差别 & 无效 & 20 例对照 \\
T11 & locate\_sigma 150 & 535.0 s（更慢） & 负优化 & 20 例对照 \\
T12 & 调参池上的 498.1 s & 独立池 538.9 / 597.0 s & \textbf{过拟合，弃用} & 两个独立池 \\
\bottomrule
\end{longtable}

% ===========================================================================
\section{复现方式}

\begin{lstlisting}[language=bash]
cd code
# 1) 快速模式（<500 s）跑一次正式测试
Q4_FAST=1 python run_http_tests.py q4 3 formal 910001
# 2) 在任意算例池上复核两种配置
python validate_fast.py 40          # 池 A 62000+、池 B 51000+，各 40 例
python validate_fast2.py 40         # 再扫一遍站位上限
# 3) 时间预算分解
python time_budget.py q4 16 30000
# 4) 全部中间实验
python cert_gap.py / pure_sweep.py / spread_scan.py / final_scan.py
\end{lstlisting}

配置定义在 \texttt{make\_data.py} 的 \texttt{P4\_FAST}；
默认仍使用精度优先的 \texttt{P4}（论文报告的那一套）。

\end{document}
